% Archivo generado por scripts/09_presentacion.py: no editar a mano
\newcommand{\periodoMeteo}{2025-01-01 a 2026-10-31}
\newcommand{\umbralParacas}{10}
\newcommand{\ventanaParacas}{12--18}
\newcommand{\mesCandidato}{2025-01}
\newcommand{\eventosCandidato}{10}
\newcommand{\criterioCandidato}{fijado en config.yaml (vientos\_paracas.mes\_candidato)}
\newcommand{\fuenteCandidato}{METAR}
\newcommand{\eventosMetar}{101}
\newcommand{\eventosEra}{0}
\newcommand{\mesMaxMetar}{2025-01 (10 días)}
\newcommand{\vientoFuerteObs}{11.0}
\newcommand{\vientoFuerteEra}{6.4}
\newcommand{\tablaValidacion}{velocidad (m/s) & 15984 & 0.20 & 1.96 & 0.59 \\
velocidad tarde 12-18 h (m/s) & 4665 & -0.46 & 2.29 & 0.36 \\
velocidad cuando obs > 10 m/s & 287 & -4.57 & 4.68 & 0.48 \\
dirección (grados) & 12399 & -25.28 & 52.61 & -- \\
temperatura (°C) & 15981 & -0.25 & 1.18 & 0.92 \\}
\newcommand{\seriePrincipal}{UHSLC}
\newcommand{\anioUhslc}{2022}
\newcommand{\qcAprobados}{1 de 16}
\newcommand{\pendienteIoc}{0.73}
\newcommand{\factorF}{0.75}
\newcommand{\claseMarea}{mixta, predominantemente semidiurna}
\newcommand{\ampMdos}{31.1}
\newcommand{\ampSdos}{9.2}
\newcommand{\ampNdos}{8.8}
\newcommand{\ampKuno}{20.0}
\newcommand{\ampOuno}{10.3}
\newcommand{\ampPuno}{6.1}
\newcommand{\rangoSicigia}{0.81}
\newcommand{\rangoCuadratura}{0.44}
\newcommand{\rangoMaximo}{1.41}
\newcommand{\rangoSicigiaObs}{1.32}
\newcommand{\rangoCuadraturaObs}{0.72}
\newcommand{\varExplicada}{93.8}
\newcommand{\rmsResiduo}{7.6}
\newcommand{\seriePrincipalNombre}{UHSLC 2022}
\newcommand{\tablaArmonicos}{M2 & 12.42 & 31.0 $\pm$ 0.1 & 337.5 $\pm$ 0.2 & 382060 \\
K1 & 23.93 & 20.0 $\pm$ 0.1 & 28.2 $\pm$ 0.3 & 94660 \\
O1 & 25.82 & 10.3 $\pm$ 0.1 & 347.3 $\pm$ 0.7 & 22187 \\
S2 & 12.00 & 9.2 $\pm$ 0.1 & 350.2 $\pm$ 0.5 & 35827 \\
N2 & 12.66 & 8.8 $\pm$ 0.1 & 297.8 $\pm$ 0.7 & 21084 \\
P1 & 24.07 & 6.1 $\pm$ 0.1 & 26.4 $\pm$ 1.0 & 10813 \\}
\newcommand{\dAmMdosIoc}{-0.9}
\newcommand{\dgMdosIoc}{4.3}
\newcommand{\areaBahia}{\pendiente}
\newcommand{\volBahia}{\pendiente}
\newcommand{\profMedia}{\pendiente}
\newcommand{\profMax}{\pendiente}
\newcommand{\dAmMdosFes}{\pendiente}
\newcommand{\dgMdosFes}{\pendiente}
